\cleardoublepage % memastikan bab baru dimulai di halaman ganjil (kanan)
\chapter{TINJAUAN PUSTAKA}
\label{chap:tinjauan-pustaka}
\section{Tinjauan Literatur}
Beberapa penelitian sebelumnya menunjukkan keberhasilan penerapan \textit{Enterprise Architecture} untuk mendukung transformasi digital sektor publik dan pertahanan, antara lain:
\begin{longtable}{|c|p{3cm}|p{3cm}|p{3cm}|p{3cm}|}
\caption{Tinjauan Literatur} \\
\hline
No & Judul & Penulis & Permasalahan & Tujuan \\
\hline
\endfirsthead

\hline
No & Judul & Penulis & Permasalahan & Tujuan \\
\hline
\endhead

1 & \textit{Implementing Cybersecurity in DefTech Sdn Bhd using DoDAF} 
& Yunnyzar Mohd Zain, Nur Azaliah Abu Bakar, Surya Sumarni Hussien (2023) \parencite{Zain2023}
& Perusahaan pertahanan (DefTech) memiliki kerangka arsitektur yang belum mengintegrasikan keamanan siber dan interoperabilitas antar divisi secara terpadu. 
& Menggunakan kerangka DoDAF untuk mendesain arsitektur “as-is” dan “to-be” di Divisi keamanan \& pertahanan perusahaan tersebut. \\
\hline
2 & \textit{Systematic Literature Review: The Role of Enterprise Architecture in Digital Transformation in the Era of the Fourth Industrial Revolution} 
& Tri Septiar Syamfithriani, Yudi Wibisono, Munir Munir, Lili Adi Wibowo, Agus Rahayu (2024) \parencite{Syamfithriani}
& Kebutuhan untuk memahami bagaimana EA mendukung digital transformation (DT) pada era Industri 4.0: teknologi baru, kerangka kerja EA yang makin banyak
& Melakukan SLR (\textit{systematic literature review}) pada literatur 2018-2022 di sektor industri untuk mengidentifikasi peranan EA dalam DT. \\
\hline
3 & \textit{Benefits of Digital Transformation and Implementation Proposition in the Defense Industry: A Literature Review}
& Mz Lerry Frendiana \& Dwi Soediantono (2022)
& Industri pertahanan menghadapi kebutuhan digitalisasi, tetapi literatur belum cukup membahas kebutuhan khusus dan keuntungan implementasi DT di sektor pertahanan.
& Mereview literatur (2015-2021) tentang digital transformation di berbagai industri dan membuat proposisi penerapan untuk industri pertahanan. \\
\hline
4 & \textit{Governance of Digital Transformation: As observed in two cases of military transformations}
& J. Kai Mattila (2022) \parencite{Mattila2022}
& Transformasi digital di organisasi militer sangat kompleks: banyak proses, sistem, manusia, teknologi—dan sering gagal realisasi.
& Menganalisis metode tata kelola (governance) yang diterapkan dalam dua studi kasus militer untuk mendukung transformasi digital. \\
\hline
5 & \textit{Enabling Open Architecture in Military Systems: A Systemic and Holistic Analysis}
& R. L. V. R. et al.(2023) \parencite{Radoman2025}
& Sistem militer menghadapi: \textit{obsolescence}, berkembangnya kebutuhan operasional, inter-operabilitas dalam SoS; banyak hambatan teknis, komersial dan organisasional.
& Memberikan analisis menyeluruh (holistik) terhadap penerapan arsitektur terbuka (\textit{Open Architecture, OA}) di sistem militer. \\
\hline
6 & \textit{Architecture Design Space Generation via Decision Pattern-Guided Department of Defense Architecture Framework Modeling}
& Zhemei Fang, Xuemeng Zhao, Fengyun Li (2024) \parencite{Fang2024}
& Kompleksitas tinggi dalam desain arsitektur sistem dan \textit{system-of-systems (SoS)}; keputusan arsitektur sering dilakukan tanpa support formal.
& Menggabungkan pola keputusan formal ke dalam EA/DoDAF untuk konteks militer; terutama terkait ruang desain arsitektur (\textit{design-space}) sistem pertahanan. \\
\hline
\end{longtable}

\subsection{Kesenjangan Penelitian (Research Gap)}
Meskipun berbagai penelitian telah membahas penerapan \textit{Enterprise Architecture} pada sektor publik dan pertahanan, 
masih terdapat kesenjangan (gap) dalam konteks pertahanan Indonesia, yaitu:
\begin{enumerate}
  \item Belum adanya model \textit{Enterprise Architecture} yang dirancang khusus untuk struktur dan sistem pertahanan nasional berbasis kebijakan Kemhan RI;
  \item Belum tersedia blueprint arsitektur pertahanan digital yang selaras dengan Arsitektur SPBE Nasional.
\end{enumerate}
Penelitian ini hadir untuk menjawab kesenjangan tersebut melalui perancangan model desain arsitektur pada Kemhan RI sebagai landasan konseptual untuk melakukan transformasi digital pertahanan nasional periode 2025–2029.

\section{Transformasi Digital}
Transformasi digital merupakan konsep strategis yang menggambarkan bagaimana organisasi memanfaatkan teknologi digital untuk menciptakan nilai baru,
meningkatkan efisiensi, dan melakukan perubahan mendasar pada proses bisnis, tata kelola, serta budaya organisasi.
Dalam literatur, transformasi digital dipandang bukan sekadar digitalisasi atau adopsi teknologi, tetapi suatu proses holistik yang mengubah cara organisasi beroperasi,
membuat keputusan, dan berinteraksi dengan pemangku kepentingan.

\begin{figure}[h] % pilihan opsi yang disarankan: t = top, b = bottom, h = here
  \caption{Digital Transformation by Verkuli} 
	\label{fig:Digital Transformation by Verkuli}
    	\includegraphics[width=0.7\textwidth]{images/DT_Verkuli.png}
\end{figure}

Menurut \textcite{Verkuli2019}, transformasi digital adalah perubahan menyeluruh berbasis teknologi yang memengaruhi strategi, organisasi, kultur, teknologi, pelanggan, dan manusia untuk menciptakan nilai dan kompetensi baru.
Definisi ini menekankan bahwa transformasi digital memiliki dimensi multiaspek yang melibatkan perubahan struktural dan kultural dalam organisasi, bukan hanya implementasi alat teknologi semata.

\begin{figure}[h] % pilihan opsi yang disarankan: t = top, b = bottom, h = here
  \caption{The Digital Transformation Canvas} 
	\label{fig:The Digital Transformation Canvas}
    	\includegraphics[width=0.7\textwidth]{images/DT_Canvas.png}
\end{figure}

Menurut \textcite{Elia2024}, menekankan bahwa transformasi digital merupakan proses strategis yang memanfaatkan teknologi digital untuk memperbarui dan mengoptimalkan tujuan organisasi,
proses internal, kolaborasi dengan mitra, pembangunan platform digital, perlindungan data, dan performa organisasi.
Transformasi digital dalam pandangan ini mencakup pembaruan pada sepuluh elemen kunci, antara lain: tujuan, proses, mitra, pembangunan kompetensi, platform, perlindungan, data dan privasi, performa, dan sumber daya manusia.
Dengan demikian, transformasi digital dipahami sebagai proses terpadu yang menghubungkan kemampuan teknologi dengan kapabilitas organisasi untuk menciptakan nilai baru.

\begin{figure}[h] % pilihan opsi yang disarankan: t = top, b = bottom, h = here
  \caption{Transformasi Digital Nasional} 
	\label{fig:Transformasi Digital Nasional}
    	\includegraphics[width=0.7\textwidth]{images/DT_Nasional.png}
\end{figure}

Di tingkat nasional, dokumen \autocite{RPJPN2025} memandang transformasi digital sebagai proses nasional yang bertujuan membangun ekosistem digital yang tangguh dan berdaulat.
Ekosistem ini dibangun melalui pengembangan infrastruktur digital, penguatan SDM digital, penyusunan regulasi pendukung, integrasi data lintas sektor, pembiayaan teknologi, serta riset dan inovasi.
Kerangka RPJPN ini memperlihatkan bahwa transformasi digital bukan hanya isu teknologi, tetapi agenda pembangunan nasional yang menuntut keterlibatan berbagai elemen, termasuk pemerintah, akademisi, industri, dan masyarakat.

Pada sektor pertahanan, transformasi digital memiliki urgensi yang lebih tinggi karena perkembangan ancaman dan dinamika perang modern sangat dipengaruhi oleh superioritas informasi, kapabilitas siber, integrasi sistem, serta kecepatan respons terhadap situasi strategis.
Hal ini menjadikan transformasi digital sebagai fondasi utama dalam membangun pertahanan yang kuat, responsif, dan adaptif terhadap perubahan lingkungan strategis global.


\section{Transformasi Digital di Sektor Pertahanan}
Transformasi digital dalam sektor pertahanan merupakan langkah strategis untuk memperkuat kapabilitas militer dan menjaga keunggulan kompetitif suatu negara dalam menghadapi ancaman kontemporer.
Berbeda dengan sektor publik pada umumnya, sektor pertahanan memiliki kompleksitas yang lebih tinggi karena mencakup pengelolaan informasi strategis, interoperabilitas sistem senjata, keamanan siber tingkat tinggi, dan dukungan terhadap operasi militer yang bersifat \textit{real-time}.
Oleh karena itu, transformasi digital di sektor pertahanan tidak hanya berfokus pada efisiensi proses internal, tetapi juga pada kesiapan tempur, kemampuan deteksi dini, dan perlindungan terhadap ancaman multidimensi, termasuk ancaman siber.

\subsection{Pendorong Transformasi Digital Pertahanan}
Beberapa faktor utama yang mendorong transformasi digital di sektor pertahanan antara lain:
\begin{enumerate}
  \item Ancaman Siber yang Semakin Kompleks.
  Ancaman siber berkembang menjadi salah satu dimensi utama perang modern.
  Negara-negara menghadapi serangan siber terhadap infrastruktur kritis, sistem militer, dan data strategis.
  Dokumen \textcite{U.S.DepartmentofDefense2019} menegaskan bahwa aktivitas siber berbahaya meningkat secara global dan setiap jaringan, sistem, serta aplikasi harus dirancang dengan prinsip “\textit{cyber first}” untuk mempertahankan postur pertahanan yang tangguh.
  \item Perubahan Karakter Perang.
  Konflik modern menekankan superioritas informasi, \textit{sensor-to-shooter integration}, \textit{command and control digital}, hingga penggunaan kecerdasan buatan dalam pengambilan keputusan.
  Oleh karena itu, platform digital menjadi fondasi sistem pertahanan yang responsif dan terintegrasi.
  \item Kebutuhan Integrasi Data dan Sistem.
  Sektor pertahanan mengelola data strategis terkait alutsista, logistik, operasi, personel, industri pertahanan, dan intelijen.
  Ketidakseragaman dan fragmentasi sistem dapat menghambat interoperabilitas, akurasi analisis, dan kecepatan pengambilan keputusan.
  \item Ekspektasi Modernisasi Organisasi dan SDM.
  Negara-negara maju, seperti Amerika Serikat, menjadikan \textit{digital workforce} sebagai salah satu pilar utama modernisasi pertahanan.
  Strategi US DoD menekankan perlunya talenta digital yang terlatih dan berkualifikasi tinggi untuk mendukung operasi pertahanan modern.
\end{enumerate}

\subsection{Praktik Transformasi Digital Pertahanan di Dunia}
\begin{enumerate}
  \item Amerika Serikat – DoD Digital Modernization Strategy.
  \begin{figure}[h] % pilihan opsi yang disarankan: t = top, b = bottom, h = here
  \caption{DoD Digital Modernization Strategy} 
	\label{fig:DoD Digital Modernization Strategy}
    	\includegraphics[width=0.7\textwidth]{images/DoD_Digital_Modernization_Strategy.png}
  \end{figure}
  DoD Digital Modernization Strategy merupakan dokumen strategis resmi Departemen Pertahanan Amerika Serikat yang menjadi landasan utama 
  dalam pelaksanaan transformasi digital di sektor pertahanan \autocite{U.S.DepartmentofDefense2019}.
  Strategi ini disusun untuk merespons perubahan karakter peperangan modern yang semakin dipengaruhi oleh perkembangan teknologi informasi, dominasi data, serta ancaman siber yang bersifat kompleks dan multidomain.
  Dalam konteks ini, keunggulan militer tidak lagi hanya ditentukan oleh kekuatan persenjataan konvensional, melainkan oleh kemampuan mengelola informasi, mengintegrasikan sistem, dan mengambil keputusan secara cepat dan akurat di seluruh domain operasi.

  DoD AS membagi transformasi digital pertahanan dalam empat pilar:
  \begin{enumerate}
    \item \textit{Innovate for competitive advantage}. 
    Inovasi adalah kunci kesiapan masa depan.
    Inovasi sangat penting untuk mempertahankan dan memperluas keunggulan dalam menghadapi persaingan dan ancaman.
    \item \textit{Optimize for efficiencies and improved capability}.
    Untuk menghadirkan kapabilitas teknologi informasi (TI) dengan efisiensi dan kinerja yang lebih tinggi, Departemen pertahanan harus mereformasi cara operasionalnya.
    Implementasi teknologi harus dipercepat.
    \item \textit{Evolve cybersecurity for resilient defense posture}.
    Aktivitas siber berbahaya terus meningkat secara global. Departemen harus mengadopsi pola pikir “Utamakan siber, selalu siber”.
    Setiap jaringan, sistem, aplikasi dan layanan perusahaan harus aman sejak awal.
    \item \textit{Cultivate talent for a digital workforce}.
    Tujuannya untuk menghasilkan tenaga kerja yang memadai,
    terdiri dari para profesional yang terlatih dan berkualifikasi tinggi untuk mendukung dan mempertahankan informasi departemen pertahanan.
  \end{enumerate}
  \item Singapura – Military Defence dalam Total Defence Strategy.
  Strategi \textit{Total Defence} Singapura merupakan pendekatan pertahanan nasional yang bersifat menyeluruh, terintegrasi, dan lintas sektor, yang pertama kali diperkenalkan pada tahun 1984.
  Strategi ini lahir dari kesadaran pemerintah Singapura bahwa keamanan nasional tidak dapat dipahami secara sempit sebagai urusan militer semata\autocite{Matthews2024}.
  Sebagai negara kota (\textit{city-state}) dengan wilayah kecil, sumber daya alam terbatas, dan masyarakat multietnis, Singapura memandang bahwa kelangsungan hidup negara sangat bergantung pada kemampuan seluruh elemen bangsa pemerintah, militer, ekonomi, dan masyarakat untuk berkontribusi secara kolektif dalam menjaga keamanan dan ketahanan nasional.

  \textit{Total Defence} dibangun di atas prinsip \textit{whole of government} dan \textit{whole of society}, di mana pertahanan diposisikan sebagai tanggung jawab bersama, bukan hanya Kementerian Pertahanan (MINDEF) dan Angkatan Bersenjata Singapura (SAF).
  Dalam perkembangannya, \textit{Total Defence} berkembang menjadi sebuah arsitektur pertahanan nasional yang terdiri dari enam pilar utama, yaitu \textit{Military Defence, Civil Defence, Economic Defence, Social Defence, Psychological Defence} dan \textit{Digital Defence}.
  Keenam pilar ini saling terhubung dan membentuk satu sistem pertahanan yang utuh.

  \begin{figure}[h] % pilihan opsi yang disarankan: t = top, b = bottom, h = here
  \caption{Singapore Total Defence} 
	\label{fig:Singapore Total Defence}
    	\includegraphics[width=0.7\textwidth]{images/Singapore_Total_Defence.png}
  \end{figure}

  Secara keseluruhan, \textit{Total Defence Strategy} menunjukkan bahwa Singapura membangun pertahanan nasionalnya melalui arsitektur komprehensif yang mengintegrasikan aspek militer, ekonomi, sosial, psikologis, dan digital dalam satu kerangka strategis.
  Pendekatan ini memungkinkan Singapura mengatasi keterbatasan fisik dan demografis melalui keunggulan teknologi, tata kelola yang kuat, serta partisipasi aktif seluruh masyarakat.
  \textit{Total Defence} bukan sekadar strategi pertahanan, melainkan sebuah model ketahanan nasional yang menempatkan keamanan dan pembangunan sebagai dua hal yang saling memperkuat.
\end{enumerate}

\section{Transformasi Digital Kemhan RI}
Transformasi digital di Kementerian Pertahanan Republik Indonesia (Kemhan RI) merupakan bagian penting dari upaya modernisasi pertahanan negara.
Seiring meningkatnya kompleksitas ancaman global, terutama pada domain siber, serta tuntutan tata kelola pemerintahan berbasis elektronik, Kemhan perlu membangun ekosistem digital yang terintegrasi, aman, dan mampu mendukung seluruh proses pertahanan secara efektif.
Transformasi ini tidak hanya menyangkut penggunaan teknologi, tetapi juga pembenahan struktur organisasi, peningkatan kapasitas SDM, serta integrasi data pertahanan untuk mendukung pengambilan keputusan yang cepat dan akurat.

\subsection{Urgensi Transformasi Digital di Kemhan}
Berikut ini merupakan beberapa alasan mendasar mengapa Kemhan perlu melakukan transformasi digital yang terarah dan menyeluruh:
\begin{enumerate}
  \item Fragmentasi sistem dan data pertahanan.
  Saat ini sistem informasi Kemhan dan matra masih berjalan secara terpisah sehingga menyebabkan kesulitan integrasi data, redundansi aplikasi, dan lemahnya interoperabilitas antar-unit.
  Hal ini berdampak pada lambatnya proses pengambilan keputusan serta kurang optimalnya perencanaan pertahanan.
  \item Meningkatnya ancaman siber.
  Ancaman siber menjadi isu strategis yang menuntut adanya penguatan sistem pertahanan siber nasional,
  termasuk perlindungan terhadap aset data dan infrastruktur digital Kemhan.
  \item Tuntutan kebijakan nasional tekait SPBE dan Satu Data Indonesia.
  \autocite{KementerianKomunikasidanInformatika2018}, dan \autocite{RepublikIndonesia2023} mengamanatkan Kemhan untuk menyusun arsitektur SPBE, peta rencana SPBE, serta layanan digital pertahanan.
  \item Arah modernisasi pertahanan dalam RPJPN 2025–2045.
  RPJPN menegaskan bahwa transformasi digital merupakan fondasi untuk menciptakan ekosistem digital pertahanan yang tangguh dan berdaulat, termasuk pembangunan infrastruktur digital, SDM digital, riset, inovasi, dan keamanan siber.
  \item Kebutuhan SDM digital pertahanan yang adaptif.
  Kemhan menghadapi tantangan keterbatasan kompetensi digital SDM sehingga modernisasi pertahanan memerlukan peningkatan literasi digital dan kemampuan manajemen data, teknologi, serta keamanan siber.
\end{enumerate}

\subsection{Tantangan Transformasi Digital Kemhan}
Tantangan utama yang harus diatasi untuk mewujudkan transformasi digital di Kemhan meliputi:
\begin{enumerate}
  \item Koordinasi lintas unit dan lintas matra yang kompleks.
  \item Kesenjangan kompetensi digital SDM.
\end{enumerate}

\section{Pertahanan Cerdas – Smart Defence}
\subsection{Definisi}
Pertahanan cerdas (\textit{Smart Defence}) merupakan konsep pertahanan modern yang menekankan pemanfaatan teknologi digital, data, dan sistem terintegrasi untuk meningkatkan efektivitas, efisiensi, serta keunggulan strategis pertahanan negara.
Konsep ini berkembang sebagai respons terhadap perubahan karakter ancaman yang semakin kompleks, multidimensi, dan berbasis teknologi informasi.
Dalam paradigma \textit{Smart Defence}, kekuatan militer tidak lagi hanya ditentukan oleh jumlah alutsista atau personel, tetapi oleh kemampuan mengelola informasi, mengintegrasikan sistem, dan mengambil keputusan secara cepat dan akurat berbasis data.

Menurut \textcite{scribd_network_centric_warfare}, \textit{Smart Defence} berakar pada pemanfaatan information superiority sebagai sumber daya strategis utama dalam peperangan modern, di mana data, jaringan, dan kecerdasan buatan menjadi komponen kritikal sistem pertahanan.
Pertahanan cerdas memandang sistem pertahanan sebagai sebuah \textit{socio-technical system} yang mengintegrasikan manusia, teknologi, proses, dan kebijakan dalam satu kesatuan arsitektur digital yang adaptif dan resilien.

Dalam konteks transformasi digital pertahanan, \textit{Smart Defence} berfungsi sebagai kerangka konseptual yang menjembatani modernisasi teknologi dengan reformasi tata kelola organisasi pertahanan.
Oleh karena itu, \textit{Smart Defence} tidak hanya berfokus pada aspek operasional militer, tetapi juga mencakup penguatan tata kelola data, keamanan siber, interoperabilitas sistem, serta pengembangan sumber daya manusia digital.

\subsection{Pergeseran Paradigma Data sebagai Weapon System}
Perkembangan teknologi digital telah mendorong terjadinya pergeseran paradigma fundamental dalam dunia pertahanan, dari \textit{platform-centric warfare} menuju \textit{data-centric warfare}.
Dalam paradigma ini, data tidak lagi diposisikan sebagai produk sampingan dari operasi militer, melainkan sebagai weapon system atau sistem senjata strategis yang memiliki nilai operasional dan strategis yang sangat tinggi.

Data yang dihasilkan dari bagian intelijen, logistik, latihan tempur, serta operasi militer menjadi dasar utama dalam proses \textit{situational awareness}, perencanaan operasi, dan pengambilan keputusan.
Keunggulan dalam mengumpulkan, mengolah, dan mendistribusikan data secara cepat dan akurat akan menghasilkan keunggulan informasi (\textit{information superiority}), yang pada akhirnya menentukan keunggulan militer di medan operasi.

Dalam konteks ini, data pertahanan tidak hanya berfungsi sebagai arsip administratif, tetapi menjadi aset strategis yang harus dikelola melalui tata kelola data (\textit{data governance}) yang kuat, standar interoperabilitas, serta mekanisme keamanan berlapis.
Tanpa arsitektur data yang terintegrasi, potensi \textit{Smart Defence} tidak dapat diwujudkan secara optimal.

\subsection{\textit{Network Centric Warfare}}
\textit{Network-Centric Warfare} (NCW) merupakan fondasi konseptual utama dari \textit{Smart Defence}.
NCW menekankan bahwa kekuatan tempur modern berasal dari kemampuan menghubungkan sensor, sistem komando dan kendali (C2), unit tempur, serta pengambil keputusan dalam satu jaringan informasi terpadu.
Melalui jaringan tersebut, informasi dapat dibagikan secara \textit{real-time}, meningkatkan \textit{shared situational awareness} dan kualitas keputusan operasional.

\begin{figure}[h] % pilihan opsi yang disarankan: t = top, b = bottom, h = here
\caption{Arsitektur NCW - Nato} 
\label{fig:Arsitektur NCW - Nato}
  \includegraphics[width=0.7\textwidth]{images/Arsitektur_NCW.png}
\end{figure}

Gambar \ref{fig:Arsitektur NCW - Nato} menunjukkan bagaimana sistem komando, sensor, dan platform tempur dihubungkan melalui jaringan komunikasi yang tersebar pada tiga lapisan utama, yaitu \textit{Ground Layer}, \textit{Airborne Layer}, dan \textit{Space Layer} \autocite{Koch2015}.
Arsitektur ini dirancang untuk mendukung konsep \textit{Network-Centric Warfare}, di mana keunggulan operasi dicapai melalui keunggulan informasi dan konektivitas jaringan.

Pada \textit{Ground Layer}, terdapat berbagai simpul komunikasi taktis seperti \textit{tactical communications node}, \textit{joint network node}, \textit{satellite tactical terminal}, \textit{super high frequency tactical satellite terminal}, serta \textit{network operations and security center}.
Lapisan ini berfungsi sebagai basis utama operasi darat, tempat data dari sensor lapangan, kendaraan tempur, dan satuan manuver dikumpulkan, diproses, serta didistribusikan.
Keberadaan \textit{network operations and security center} menunjukkan bahwa selain konektivitas, aspek keamanan jaringan dan pengendalian operasi menjadi elemen kunci dalam NCW.

Di atasnya, \textit{Airborne Layer} berperan sebagai penguat jangkauan (\textit{reach}) dan penghubung fleksibel antara darat dan ruang angkasa.
Platform seperti UAV jarak jauh (misalnya \textit{Global Hawk} dan \textit{extended-range unmanned aircraft systems}) berfungsi tidak hanya sebagai sensor intelijen, tetapi juga sebagai relay komunikasi udara.
Lapisan ini sangat penting untuk mengatasi keterbatasan geografis, gangguan medan, maupun kerusakan infrastruktur darat, sekaligus meningkatkan ketahanan jaringan dalam kondisi pertempuran.

Sementara itu, \textit{Space Layer} menyediakan kemampuan \textit{reach \& reachback} secara global melalui berbagai sistem satelit, baik \textit{military SATCOM}, \textit{advanced extremely high frequency (EHF)}, \textit{wideband global SATCOM}, maupun satelit komersial.
Lapisan ini memungkinkan komunikasi jarak jauh, pengendalian operasi lintas wilayah, serta koneksi antara satuan tempur di medan operasi dengan pusat komando strategis.
Keberadaan berbagai jenis satelit menunjukkan strategi redundansi dan diversifikasi jalur komunikasi, guna mengurangi ketergantungan pada satu sistem tunggal.

Garis-garis penghubung antarobjek dalam gambar menegaskan bahwa arsitektur NCW bersifat \textit{many-to-many connectivity}, bukan hubungan linier.
Artinya, setiap simpul dapat saling bertukar data melalui berbagai jalur alternatif, sehingga jaringan tetap dapat berfungsi meskipun sebagian node mengalami gangguan atau diserang.
Konsep ini sangat relevan dalam menghadapi ancaman \textit{Anti-Access/Area Denial (A2AD)}, di mana lawan berupaya melumpuhkan komunikasi dan sistem komando.

Secara keseluruhan, Gambar \ref{fig:Arsitektur NCW - Nato} merepresentasikan ekosistem komunikasi pertahanan modern yang terintegrasi lintas domain darat, udara, dan antariksa.
Arsitektur tersebut menegaskan bahwa keberhasilan operasi militer modern tidak hanya ditentukan oleh kekuatan senjata, tetapi juga oleh ketahanan, fleksibilitas, dan keamanan jaringan informasi sebagai fondasi utama pertahanan cerdas.

\subsection{\textit{Join All-Domain Operation}}
Seiring perkembangan NCW, konsep pertahanan modern berevolusi menuju \textit{Joint All-Domain Operations (JADO)}, yaitu pendekatan operasi militer yang mengintegrasikan seluruh domain operasi darat, laut, udara dan siber dalam satu kerangka komando dan kendali terpadu.
JADO menuntut integrasi lintas domain secara simultan untuk menciptakan keunggulan operasional yang berkelanjutan.

US Department of Defense (2022) menerbitkan strategi \textit{Joint All Domain Command and Control (JADC2)} sebagai arah strategis modernisasi sistem Komando dan Kendali (\textit{Command and Control/C2}) dalam menghadapi dinamika lingkungan keamanan global yang semakin kompleks dan kompetitif \autocite{USDoD2022}.
Strategi ini lahir dari kesadaran bahwa keunggulan militer tidak lagi semata ditentukan oleh kekuatan fisik dan alutsista, melainkan oleh kemampuan menguasai informasi serta kecepatan dan ketepatan pengambilan keputusan dalam operasi lintas domain.

Strategi JADC2 dibangun di atas tiga fungsi utama C2, yaitu \textit{sense}, \textit{make sense}, dan \textit{act}.
Fungsi \textit{sense} merujuk pada kemampuan mengumpulkan, mengintegrasikan, dan membagikan data dari seluruh domain operasi darat, laut, udara, siber, dan antariksa serta dari berbagai sumber, termasuk sensor militer, intelijen, dan mitra operasi.
Integrasi ini dilakukan melalui konsep \textit{federated data fabric} yang memungkinkan pertukaran data lintas sistem secara aman dan real-time.

Selanjutnya, fungsi \textit{make sense} berfokus pada proses analisis dan pemaknaan data untuk memahami kondisi operasional, memprediksi tindakan lawan, serta mengevaluasi posisi dan kemampuan pasukan sendiri dan mitra.
Pada tahap ini, data yang telah dikumpulkan diolah menjadi informasi dan pengetahuan melalui pemanfaatan teknologi analitik canggih, termasuk \textit{Artificial Intelligence (AI)} dan \textit{Machine Learning (ML)}.
Pemanfaatan teknologi tersebut ditujukan untuk mempercepat siklus pengambilan keputusan komandan sehingga mampu bertindak lebih cepat dibandingkan lawan, bahkan dalam kondisi C2 yang terganggu atau terbatas.

Fungsi ketiga, yaitu \textit{act}, berkaitan dengan kemampuan mengambil keputusan dan mendistribusikannya secara cepat, akurat, dan aman ke seluruh unsur \textit{Joint Force} dan mitra operasi.
Fungsi ini menekankan pentingnya kombinasi antara pengambilan keputusan manusia dan dukungan sistem digital, serta penerapan prinsip \textit{mission command}.
Melalui pendekatan ini, satuan bawah tetap memiliki kewenangan dan pemahaman terhadap niat operasi pimpinan sehingga mampu bertindak secara mandiri ketika komunikasi terputus atau situasi menuntut respons segera.

\begin{figure}[h] % pilihan opsi yang disarankan: t = top, b = bottom, h = here
\caption{Arsitektur JADC2} 
\label{fig:Arsitektur JADC2}
  \includegraphics[width=0.7\textwidth]{images/Arsitektur_JADC2.png}
\end{figure}

Secara keseluruhan, JADC2 menegaskan bahwa keberhasilan operasi militer modern sangat bergantung pada kemampuan integrasi informasi dan pengambilan keputusan lintas domain secara cepat dan adaptif.

\subsection{Pilar Utama Modernisasi Pertahanan Cerdas Kemhan RI}
\begin{enumerate}
  \item Pilar Data \& Informasi Pertahanan.
  Pilar Data \& Informasi Pertahanan merupakan inti dari konsep pertahanan cerdas.
  Pilar ini menempatkan data sebagai aset strategis negara yang menjadi dasar keunggulan informasi (\textit{information superiority}).
  Data pertahanan mencakup data intelijen, operasi, logistik, alutsista, personel, hingga data ancaman siber yang berasal dari berbagai sumber dan sensor.
  \item Pilar Teknologi Digital \& Sistem Cerdas.
  Pilar ini berfokus pada pemanfaatan teknologi digital mutakhir sebagai enabler pertahanan cerdas.
  Teknologi seperti Artificial Intelligence (AI), Big Data Analytics, Cloud \& Edge Computing, Internet of Military Things (IoMT), serta sistem otonom digunakan untuk mengolah data pertahanan menjadi informasi dan rekomendasi keputusan yang bernilai strategis.
  \item Pilar Integrasi Sistem \& Interoperabilitas.
  Pilar Integrasi Sistem \& Interoperabilitas menekankan pentingnya keterhubungan lintas sistem, organisasi, dan matra.
  Pertahanan cerdas mensyaratkan integrasi antara Kemhan, TNI AD, AL, AU, serta pemangku kepentingan pertahanan lainnya dalam satu arsitektur pertahanan terpadu.
  \item Pilar SDM Pertahanan Digital.
  Transformasi menuju pertahanan cerdas menuntut sumber daya manusia pertahanan yang unggul, adaptif, dan memiliki literasi digital tinggi.
  Pilar ini mencakup pengembangan kompetensi personel di bidang data analytics, AI, keamanan siber, dan sistem informasi pertahanan, serta pembentukan budaya kerja berbasis data (data-driven culture).
  \item Pilar Tata Kelola, Kebijakan \& Keamanan Siber.
  Pilar ini berperan sebagai pengendali dan pengaman seluruh ekosistem pertahanan cerdas.
  Tata kelola yang baik memastikan transformasi digital pertahanan berjalan selaras dengan regulasi, kebijakan nasional, dan kepentingan strategis negara.

  Keamanan siber bukan hanya isu teknis, tetapi bagian dari pengambilan keputusan pertahanan negara yang berkaitan dengan otoritas, klasifikasi informasi, serta respons terhadap ancaman siber sebagai bentuk peperangan modern.

  \item Pilar Kolaborasi \& Inovasi Industri Pertahanan.
  Pilar Kolaborasi \& Inovasi menekankan pentingnya sinergi antara Kemhan, industri pertahanan nasional, perguruan tinggi, dan lembaga riset.
  Kolaborasi ini bertujuan mempercepat inovasi teknologi pertahanan, mendorong transfer teknologi, serta memperkuat kemandirian industri pertahanan nasional.

  Dalam konteks pertahanan cerdas, inovasi tidak hanya berfokus pada alutsista, tetapi juga pada platform digital, sistem data, dan solusi cerdas.
  Pilar ini memastikan bahwa pertahanan cerdas Kemhan bersifat adaptif, berkelanjutan, dan tidak bergantung sepenuhnya pada teknologi luar negeri.
\end{enumerate}

\section{Framework \textit{Enterprise Architecture}}
Transformasi digital pada sektor pertahanan memerlukan kerangka \textit{Enterprise Architecture} (EA) yang mampu mengintegrasikan aspek organisasi, proses bisnis, data, aplikasi, teknologi, keamanan siber, dan interoperabilitas lintas domain.
Berbagai \textit{framework} EA telah dikembangkan dan diterapkan pada organisasi pemerintahan maupun pertahanan, di antaranya TOGAF, DoDAF, MODAF dan NAF.
Masing-masing \textit{framework} memiliki karakteristik, keunggulan, dan keterbatasan yang berbeda.
Oleh karena itu, dilakukan analisis untuk menentukan \textit{framework} yang paling sesuai sebagai dasar pengembangan arsitektur transformasi digital pada Kementerian Pertahanan Republik Indonesia.

\subsection{DoDAF}
\textit{Department of Defense Architecture Framework} (DoDAF) dikembangkan oleh Departemen Pertahanan Amerika Serikat untuk mendukung pengambilan keputusan melalui berbagai informasi.
Kerangka ini berevolusi dari \textit{C4ISR Architecture Framework} versi 1.0 yang dirilis pada tahun 1996 \autocite{C4ISRAWG1998}, melewati DoDAF versi 1.0 (2004) dan 1.5 (2007) \autocite{DoDAF2007v15}, hingga DoDAF versi 2.0 \autocite{DoDAF2010v20} yang disetujui pada Mei 2009 dan versi 2.02 \autocite{DoDAF2011v202} yang menjadi iterasi terkini hingga saat ini.
Pergeseran paling fundamental pada versi 2.0 adalah perubahan dari pendekatan berorientasi produk menjadi pendekatan yang bersifat \textit{data-centric} alih-alih berpusat pada produk atau layanan, melalui penggunaan \textit{DoDAF Meta-model} (DM2) sebagai dasar pengembangan arsitektur.

Secara struktural, DoDAF 2.02 memuat delapan \textit{viewpoint} dan 52 model, yang mencakup \textit{All Viewpoint} (AV), \textit{Capability Viewpoint} (CV), \textit{Data and Information Viewpoint} (DIV), \textit{Operational Viewpoint} (OV), \textit{Project Viewpoint} (PV), \textit{Services Viewpoint} (SvcV), \textit{Standards Viewpoint} (StdV), dan \textit{Systems Viewpoint} (SV).
Kemampuan utama DoDAF yaitu mengorganisasi arsitektur \textit{systems-of-systems} yang kompleks dan terdistribusi, sebagaimana dibutuhkan dalam operasi berbasis jaringan (\textit{Network Centric Operations}).
Meskipun demikian, DoDAF pada dasarnya merupakan kerangka \textit{deskripsi} arsitektur yang menekankan apa yang harus didokumentasikan, dan tidak menyediakan metodologi pengembangan bertahap (\textit{step-by-step}).

\subsection{NAF}
\textit{NATO Architecture Framework} (NAF) merupakan kerangka arsitektur yang dikembangkan oleh Pakta Pertahanan Atlantik Utara (NATO) untuk meningkatkan interoperabilitas antar negara anggota.
\textit{Framework} ini sangat kuat dalam mendukung operasi gabungan dan integrasi sistem pertahanan lintas organisasi, namun kompleksitas implementasinya relatif tinggi dan belum sepenuhnya sesuai dengan kebutuhan tata kelola pemerintahan Indonesia.

Versi terkini, NAF versi 4.1 (NAFv4) \autocite{NAFv4}, diterbitkan oleh \textit{Architecture Capability Team} (ACaT) di bawah \textit{NATO Consultation, Command and Control Board} (C3B).
NAFv4 disusun atas beberapa komponen utama, yaitu metodologi pengembangan arsitektur, \textit{viewpoint}, serta penerapan \textit{meta-model} komersial yang sesuai dengan kebijakan NATO.
Karakteristik khas NAFv4 adalah pengorganisasian \textit{viewpoint} dalam bentuk struktur ``grid'' dua dimensi yang merepresentasikan berbagai ``isu yang menjadi perhatian'' (\textit{issues of concern}) pada baris dan ``aspek yang menjadi perhatian'' (\textit{aspects of concern}) pada kolom.
NAF dirancang untuk selaras dengan referensi arsitektur yang diterbitkan oleh badan standar internasional seperti ISO dan IEEE, sehingga sangat menonjol pada aspek interoperabilitas multinasional.

\subsection{MODAF}
\textit{Ministry of Defence Architecture Framework} (MODAF) merupakan kerangka arsitektur \textit{enterprise} yang dikembangkan oleh Kementerian Pertahanan Inggris untuk mendukung perencanaan pertahanan dan manajemen perubahan \autocite{Weilkiens2022}.
MODAF merupakan adaptasi DoDAF yang masih berorientasi pada kebutuhan organisasi pertahanan Inggris sehingga memerlukan penyesuaian signifikan untuk diterapkan pada konteks Indonesia.

MODAF menambahkan dua \textit{viewpoint} baru, yaitu \textit{Strategic Viewpoint} dan \textit{Acquisition Viewpoint}.
Secara keseluruhan, MODAF mendefinisikan tujuh \textit{viewpoint}, yaitu \textit{All Views} (AV), \textit{Strategic Views} (StV), \textit{Operational Views} (OV), \textit{Service-Oriented Views} (SOV), \textit{Systems Views} (SV), \textit{Acquisition Views} (AcV), dan \textit{Technical Views} (TV).
Penambahan \textit{viewpoint} strategis dan akuisisi ditujukan untuk menjawab kebutuhan \textit{Capability Managers} dan \textit{Acquisition Managers} dalam menganalisis evolusi kapabilitas serta keterkaitan antarprogram.

\subsection{TOGAF}
\autocite{togaf2020} merupakan salah satu kerangka kerja (\textit{enterprise architecture}) yang paling banyak digunakan secara internasional untuk membantu organisasi dalam merancang, mengelola, dan mengimplementasikan arsitektur sistem secara terstruktur.
TOGAF menyediakan metodologi yang komprehensif, standar, serta seperangkat prinsip yang memungkinkan organisasi mengelola perubahan secara efektif dan memastikan bahwa proses transformasi berjalan selaras dengan tujuan strategis organisasi.

Inti dari TOGAF adalah \textit{Architecture Development Method} (ADM), yaitu sebuah siklus metodologis yang memberikan panduan langkah demi langkah dalam merancang arsitektur organisasi.
Siklus ADM terdiri dari beberapa fase yang mencakup perencanaan awal, pengembangan arsitektur bisnis, data, aplikasi, teknologi, hingga perencanaan implementasi dan tata kelola.
Dalam penelitian ini, fokus diberikan pada fase A hingga D, yang digunakan untuk merancang arsitektur transformasi digital di lingkungan Kementerian Pertahanan RI.

Fase-fase utama yang relevan dalam penelitian ini meliputi:
\begin{enumerate}
  \item Phase A – Architecture Vision.
  Pada fase ini ditetapkan visi arsitektur, ruang lingkup, dan kebutuhan strategis yang ingin dicapai.
  Dalam konteks Kemhan RI, fase ini berfungsi untuk merumuskan tujuan transformasi digital.
  \item Phase B – Business Architecture.
  Fase ini menganalisis proses bisnis yang berjalan saat ini dan merancang proses bisnis target yang mendukung transformasi digital.
  \item Phase C – Information Systems Architecture (Data \& Application Architecture).
  Pada fase ini dilakukan perancangan arsitektur data dan aplikasi yang mendukung transformasi digital.
  Analisis meliputi kebutuhan data kompetensi, sistem manajemen pelatihan, integrasi data pelatihan, serta aplikasi pendukung berbasis digital yang dapat digunakan oleh Kemhan.
  \item Phase D – Technology Architecture.
  Fase ini merancang kebutuhan teknologi yang diperlukan untuk mendukung implementasi arsitektur transformasi digital.
  Pada konteks Kemhan, hal ini mencakup kebutuhan infrastruktur seperti platform teknologi cloud, sistem keamanan siber, dan perangkat integrasi data.
\end{enumerate}

Penggunaan TOGAF dalam penelitian ini memberikan manfaat signifikan, yakni memastikan bahwa model arsitektur transformasi digital Kemhan RI yang dihasilkan berfokus pada proses bisnis, kebutuhan data, sistem aplikasi pendukung, dan infrastruktur teknologi yang terintegrasi.
Kerangka ini membantu organisasi menghindari pendekatan yang parsial dan memastikan bahwa transformasi digital dilakukan secara menyeluruh serta berkelanjutan.

Selain itu, pemanfaatan TOGAF ADM mendukung konsistensi dengan praktik terbaik dunia dalam pembangunan arsitektur organisasi, termasuk pada sektor pertahanan.
Dengan struktur metodologis yang jelas, TOGAF memungkinkan Kemhan RI untuk merancang arsitektur transformasi digital yang dapat diimplementasikan, dipantau, dan dievaluasi secara berkelanjutan.

\begin{figure}[h] % pilihan opsi yang disarankan: t = top, b = bottom, h = here
\caption{TOGAF Framework} 
\label{fig:TOGAF Framework}
  \includegraphics[width=0.7\textwidth]{images/TOGAF_Framework.png}
\end{figure}

\subsection{Perbandingan \textit{Framework}}
Perbandingan keempat kerangka EA tersebut disajikan pada Tabel \ref{tab:perbandingan-framework}.

{\footnotesize
\begin{longtable}{|p{2.1cm}|p{2.5cm}|p{2.5cm}|p{2.5cm}|p{2.5cm}|}
\caption{Perbandingan antar \textit{framework}} \label{tab:perbandingan-framework} \\
\hline
\textbf{Aspek} & \textbf{DoDAF} & \textbf{NAF} & \textbf{MODAF} & \textbf{TOGAF} \\
\hline
\endfirsthead

\hline
\textbf{Aspek} & \textbf{DoDAF} & \textbf{NAF} & \textbf{MODAF} & \textbf{TOGAF} \\
\hline
\endhead

Pengembang & Departemen Pertahanan AS (DoD) & NATO (\textit{Architecture Capability Team}, C3B) & Kementerian Pertahanan Inggris (UK MoD) & The Open Group \\
\hline
Versi \& Status & 2.02 (2010), aktif & v4 (2018), aktif & v1.2, ditarik (digantikan NAF v4) & 10th Edition (2022), aktif \& berkembang \\
\hline
Orientasi utama & Dokumentasi arsitektur \textit{systems-of-systems} \& C4ISR & Standardisasi \& interoperabilitas multinasional & Perencanaan kapabilitas \& akuisisi pertahanan & Keselarasan bisnis--TI lintas sektor \\
\hline
Tipe kerangka & \textit{Architecture description framework} & Deskripsi + metodologi (sejak v4) & \textit{Architecture description framework} & \textbf{Metodologi pengembangan} (proses) + \textit{content framework} \\
\hline
Komponen/ Struktur & 8 \textit{viewpoint}, 52 model, meta-model DM2 & \textit{Grid viewpoint} (\textit{issues} $\times$ \textit{aspects of concern}), meta-model komersial & 7 \textit{viewpoint} (AV, StV, OV, SOV, SV, AcV, TV) & Siklus ADM (Fase A--H), 4 domain arsitektur, \textit{Enterprise Continuum}, \textit{Architecture Repository} \\
\hline
Metode Pengembangan & Tidak preskriptif; berpusat pada data \& \textit{viewpoint} & \textit{Architecting stages} (tidak ketat sekuensial) & Tidak preskriptif; turunan DoDAF & \textbf{ADM iteratif \& preskriptif} dengan \textit{Requirements Management} di pusat \\
\hline
Cakupan Sektor & Spesifik militer (AS) & Militer (NATO) \& bisnis & Spesifik militer (Inggris) & Netral industri / lintas sektor \\
\hline
Penanganan keamanan & Akar C4ISR kuat; \textit{security markings} di PES & Fokus interoperabilitas lintas matra/negara & Kuat pada kapabilitas \& akuisisi & Tidak eksplisit; bersifat \textit{cross-cutting} (perlu adaptasi) \\
\hline
Fleksibilitas \& Kustomisasi & Terikat konvensi DoD & Selaras ISO/IEC/IEEE 42010; kompatibel mundur & Menambah StV \& AcV; kini \textit{legacy} & \textbf{Sangat tinggi} --- ADM dirancang untuk di-\textit{tailor} \\
\hline
Adopsi \& ketersediaan & Wajib bagi akuisisi sistem DoD & Standar NATO \& negara mitra & \textit{Legacy}, tidak lagi dipelihara & Adopsi global terluas; ekosistem \& sertifikasi matang; gratis non-komersial \\
\hline
\end{longtable}
}

Hasil perbandingan pada Tabel \ref{tab:perbandingan-framework} memperlihatkan adanya perbedaan mendasar dalam karakter keempat kerangka tersebut.
DoDAF, NAF, dan MODAF pada dasarnya merupakan \textit{architecture description framework}, yaitu kerangka yang menekankan apa yang harus didokumentasikan melalui sekumpulan \textit{viewpoint} dan model baku.
Ketiganya lahir dari akar yang sama, yaitu \textit{C4ISR Architecture Framework}, dan dirancang secara spesifik untuk konteks militer suatu negara atau aliansi.
Sebaliknya, TOGAF merupakan \textit{development methodology} yang menekankan \textit{bagaimana} arsitektur dibangun secara bertahap melalui \textit{Architecture Development Method} (ADM), serta bersifat netral industri dan dapat diterapkan lintas sektor.
Perbedaan inilah yang menjadi pertimbangan utama dalam memilih TOGAF sebagai dasar pengembangan kerangka pada penelitian ini.

Keunggulan metodologis TOGAF terletak pada siklus ADM yang sistematis, iteratif, dan preskriptif.
ADM memberikan panduan langkah demi langkah mulai dari \textit{Preliminary Phase}, perumusan visi, pengembangan arsitektur bisnis, data, aplikasi, dan teknologi, hingga perencanaan implementasi dan tata kelola perubahan, dengan \textit{Requirements Management} yang konsisten berada di pusat siklus.
Dengan demikian, TOGAF tidak hanya menghasilkan dokumentasi arsitektur, tetapi juga mengarahkan proses transformasi secara menyeluruh---suatu kemampuan yang tidak dimiliki secara eksplisit oleh DoDAF dan MODAF yang bersifat deskriptif.
Kemampuan ini sangat relevan dengan kebutuhan Kementerian Pertahanan RI yang tidak hanya memerlukan peta arsitektur, melainkan metode terstruktur untuk melaksanakan transformasi digital.

Keunggulan kedua adalah fleksibilitas dan kemampuan kustomisasi.
ADM secara eksplisit dirancang untuk di-\textit{tailor} sesuai karakteristik organisasi, sehingga membuka ruang adaptasi yang luas tanpa kehilangan kerangka metodologis intinya.
Sifat ini menjadikan TOGAF sebagai fondasi yang ideal untuk dikembangkan menjadi \textit{TOGAF Pertahanan Indonesia}.
Selain itu, TOGAF merupakan kerangka arsitektur \textit{enterprise} dengan adopsi terluas secara global serta tersedia secara terbuka untuk penggunaan non-komersial dan akademik.
Karakter netral-sektor TOGAF juga memudahkan penyelarasan dengan kebijakan nasional seperti SPBE, Satu Data Indonesia, dan RPJPN 2025--2045, yang menuntut keselarasan tata kelola lintas instansi pemerintah, bukan hanya kebutuhan militer.

Meskipun demikian, perbandingan ini juga menegaskan satu keterbatasan TOGAF dalam konteks pertahanan, yaitu tidak adanya domain yang secara khusus menangani kapabilitas operasional militer seperti C5ISR serta penanganan keamanan siber yang masih bersifat \textit{cross-cutting}.
Justru pada titik inilah kekuatan DoDAF, NAF, dan MODAF menjadi relevan: ketiganya unggul dalam merepresentasikan domain komando, kendali, dan interoperabilitas yang berakar pada C4ISR/C5ISR.
Oleh karena itu, strategi yang ditempuh dalam penelitian ini adalah menggunakan kekuatan metodologis ADM TOGAF sebagai kerangka utama dan menanamkan kebutuhan domain pertahanan yang menjadi keunggulan ketiga kerangka militer tersebut, khususnya melalui penambahan fase \textit{C5ISR Architecture}.




























